\section{Phase Boundary Theorem}
\label{sec:phase_boundary}

This section presents the central result: a characterized boundary in
the parameter space of the coupled $(P, \Wm)$ system that separates
self-correcting dynamics from absorbing (or metastable) dynamics.

\begin{theorem}[Phase Boundary for Compound Feedback Loops]
\label{thm:phase_boundary}
Consider the coupled $(P, \Wm)$ dynamical system under locally rational
transitions (Definition~\ref{def:transition}), with
Assumptions~C1 and~C2$'$ (Section~\ref{sec:lyapunov}).  Define:
\begin{itemize}
\item $\rS = \inf_t \alpha(t) \cdot c_{\mathrm{S}}$: a lower bound on
  the exogenous correction rate (from Assumption~C1's first term),
\item $c_{\mathrm{V}}$, $\nu$: the endogenous correction constant and
  $\volL$-sensitivity (from Assumption~C1's second term; set
  $c_{\mathrm{V}} = 0$ for the purely exogenous model),
\item $\rW = \sup_t \beta(t) \cdot c_{\mathrm{D}}$: an upper bound on
  the error-proportional degradation rate per subsumption step,
\item $\dmax$: the maximum exogenous drift rate per dimension (from
  C2$'$; a global constant),
\item $\rsub \in [0,1]$: the time-averaged subsumption frequency over
  the trajectory window,
\item $\rhomincross = \inf_{t \in [t_0, t_0+T]} \min_k
  \rhocross_k(P_t)$: the infimum of the minimum cross-substrate
  redundancy over the window
  (Definition~\ref{def:cross_redundancy}),
\item $B(t_0, T)$: the cumulative error budget over window
  $[t_0, t_0 + T]$ (Equation~\ref{eq:error_budget}).
\end{itemize}

\noindent The simplified sufficient conditions below use these
worst-case bounds.  The full trajectory-level statement (Part~a) uses
the uniform contraction condition on the per-step multiplicative
factors (Equation~\ref{eq:uniform_contraction}), which the simplified
condition implies in the linearized regime.  \emph{Note:} the explicit
neighborhood radius (Equation~\ref{eq:neighborhood}) is derived from
worst-case instantaneous bounds $\bar{a}$, $\bar{b}$ via a linearized
approximation ($\ln a_t \approx a_t - 1$ when $|a_t - 1| \ll 1$).
The general convergence to $\mu \cdot \bar{b}/(1-\bar{a})$ holds under
the uniform contraction condition without linearization ($\mu$ is the
overshoot constant defined in the appendix); the substitution of named
parameters into $\bar{a}$ and $\bar{b}$, and the drop of $\mu$ to 1
at linearization order, require the linearized regime
(Remark~\ref{rem:linearized}).

\medskip
\noindent\textbf{Part~(a): Self-Correcting Basin.}
If the uniform contraction condition holds (there exist constants
$\bar{a} \in (0,1)$ and $\bar{b} \geq 0$ such that, for all
sufficiently long windows $[t_0, t_0+T]$, the time-averaged
log-contraction rate of the per-step multiplicative factor
$a_t$ satisfies $\frac{1}{T}\sum \ln a_t \leq \ln \bar{a} < 0$ and
the time-averaged additive drift satisfies
$\frac{1}{T}\sum b_t \leq \bar{b}$; see Equation~\ref{eq:uniform_contraction}
for notation), then $\Lyap(\Wm_t)$ converges to a neighborhood of zero
whose radius is bounded by $\mu \cdot \bar{b}/(1 - \bar{a})$ in
general (with overshoot constant $\mu$; see proof).  In the
linearized regime (Remark~\ref{rem:linearized}), this evaluates to
\begin{equation}
\label{eq:neighborhood}
\Lyap_\infty \leq \frac{\rsub \cdot \dmax \cdot
  \overline{\Delta\rho}}{\rS \cdot \rhomincross
  + \alphamin \cdot c_{\mathrm{V}} \cdot \underline{\nu} - \rW \cdot \rsub}
\end{equation}
where $\overline{\Delta\rho}$ denotes the time-averaged redundancy
loss per subsumption step.  (We write $\overline{\Delta\rho}$
throughout; in the linearized simplification below,
$\Delta\rho_{\mathrm{avg}}$ is used synonymously when the subscript
form improves readability.)  This neighborhood shrinks as subsumption ceases
($\overline{\Delta\rho} \to 0$).  Once $\Lyap(\Wm_t) <
\epsilon_{\mathrm{safe}}$ (defined as the error threshold below which
the locally rational policy agrees with the true $\Vdisc$-optimal
policy), subsumption is correctly evaluated against the true $\volL$
dynamics and the anti-monopolar result~\citep[Proposition~6]{lasser2026horizon} applies.

The feedback loop self-corrects provided the neighborhood radius
$\Lyap_\infty < \epsilon_{\mathrm{safe}}$, which imposes an additional
constraint on the steady-state subsumption rate.  \emph{Note:}
$\epsilon_{\mathrm{safe}} > 0$ requires that the optimal action at
$\Wm^*$ has a strict gap ($G_j(\Wm^*) > 0$ for all preserve-optimal
agents).  This holds generically for finite action sets (see proof).
A computable lower bound using action-gap gradient norms
(Cauchy--Schwarz in the dual $w^{-1}$-weighted norm, induced by the
$w$-weighted metric on $\Wm$-space) is derived there.  The
easier-to-compute per-coordinate estimator $\min_k w_k \cdot
\epsilon_{\mathrm{crit},k}^2$ rules out single-coordinate flips but
may overestimate $\epsilon_{\mathrm{safe}}$ when action gaps are
sensitive to multiple coordinates simultaneously.

\medskip
\noindent\textbf{Simplified sufficient condition.}  In the linearized
regime (Remark~\ref{rem:linearized}):
\begin{equation}
\label{eq:self_correcting}
\rS \cdot \rhomincross + \alphamin \cdot c_{\mathrm{V}} \cdot \underline{\nu}
  > \rW \cdot \rsub
  + \dmax \cdot \frac{\Delta\rho_{\mathrm{avg}}}{\Lyap_{\mathrm{avg}}}
\end{equation}
where $\alphamin = \inf_t \alpha(t)$ and
$\underline{\nu} = \inf_t \nu(P_t, \Wm_t)$.  (The per-step bound from
Proposition~\ref{prop:error_dynamics} carries the time-varying
$\alpha(t) \cdot c_{\mathrm{V}} \cdot \nu(P_t, \Wm_t)$; the
simplified condition uses the inf bounds $\alphamin$ and
$\underline{\nu}$.)  The left-hand side has two
terms: the exogenous correction $\rS \cdot \rhomincross$ (proportional
to channel strength, where $\rS = \alphamin \cdot c_{\mathrm{S}}$
already absorbs the learning-rate floor) and the endogenous incentive
$\alphamin \cdot c_{\mathrm{V}}
\cdot \underline{\nu}$ (the actor's own $\volL$-gradient toward better
calibration, scaled by the same learning-rate floor).  The endogenous term
widens the self-correcting basin by supplementing the observation
channels.  Setting $c_{\mathrm{V}} = 0$ recovers the purely exogenous
condition.

The first term on the right is the error-proportional degradation cost;
the second is the drift-exposure cost from lost channels.  When
subsumption removes channels covering well-calibrated dimensions
($\Lyap_{\mathrm{avg}}$ is low but $\Delta\rho_{\mathrm{avg}}$ is
nonzero), the second term dominates and the condition is harder to
satisfy, correctly diagnosing the ``removing a channel that was doing
all the correction work'' failure mode.

\begin{remark}[Stabilizing and destabilizing cascades]
\label{rem:stabilizing_cascade}
Part~(a) describes a \emph{stabilizing cascade}: as the world model
improves ($\Lyap$ decreases), the actor makes better decisions, which
preserve more cooperative capabilities, which preserve more channels,
which provide more correction signal.  Part~(b), below, describes the
mirror image: a \emph{destabilizing cascade} where channel loss
degrades the model, producing worse decisions that lose more channels.
The phase boundary separates these two regimes.  The endogenous
$\alphamin \cdot c_{\mathrm{V}} \cdot \nu$ term from Assumption~C1 tips the balance
toward the stabilizing side by supplementing the channel-driven
correction with the actor's own $\volL$-gradient incentive.

The two cascades are asymmetric.  Degradation happens in discrete
jumps: subsumption removes an agent and its channels at a single step.
Recovery is continuous: calibration improves gradually through
accumulated observations.  The system degrades in steps but recovers
smoothly.  The destabilizing cascade therefore has discrete tipping
points (the moment $\rhocross_k$ hits zero on a cascade dimension),
while the stabilizing cascade has no corresponding threshold, only a
smooth approach to full calibration.
\end{remark}

\medskip
\noindent\textbf{Part~(b): Absorbing Basin.}
Under Assumption~B1 (no endogenous correction on blind dimensions;
Assumption~\ref{ass:B1}) and the multi-dimensional monotonicity
condition (ii$'$) from the Cascade Lemma~(\ref{lem:cascade}; needed
for propagation beyond the initial trigger):

If there exists a trajectory segment $[t_0, t_0 + T]$ during which
\begin{enumerate}
\item[(i)] $B(t_0, T) > 0$ (cumulative degradation exceeds cumulative
  correction),
\item[(ii)] by $t_0 + T$, $\rhocross_k(P_{t_0+T})$ has been reduced to $0$
  on at least one cascade dimension $k$ with $w_k > 0$, and
\item[(iii)] the accumulated error on $k$ exceeds the cascade
  threshold: $\epsilon_k(t_0 + T) > \epsilon_{\mathrm{crit},k}$
  (Lemma~\ref{lem:cascade}; the minimum error on dimension $k$ that
  flips any preserve/subsume action),
\end{enumerate}
then the following cascade occurs.  By Assumption~B1a,
$\epsilon_k(t)$ is non-decreasing for $t > t_0 + T$ (the blinded
dimension cannot self-correct).  If the environment is
\emph{non-stationary} on dimension $k$---the true value $\Wm^*[k]$
continues to evolve while the agent's estimate $\Wm_t[k]$ remains
frozen---then $\epsilon_k$ grows and condition~(iii) is sustained.
(Assumption~C2$'$ bounds the \emph{rate} of drift via $\dmax$ but does
not guarantee drift occurs; the cascade requires the structural
condition that blinded dimensions are not environmentally stationary.
In the multi-agent setting where subsumption eliminates active
participants, non-stationarity is the generic case: the remaining
agents' actions change the true state on dimensions the eliminated
agent's channels were tracking.)

The Cascade Lemma (Appendix~\ref{app:cascade}) propagates the failure:
the resulting subsumptions drive $\rhocross_j \to 0$ on additional
dimensions $j \neq k$ (provided the agent-channel graph is
connected, so each subsumption removes channels covering at least one
new dimension).  Each newly blinded dimension $j'$ has
$\epsilon_{j'}$ non-decreasing (by B1a); if environmental
non-stationarity causes $\epsilon_{j'}$ to exceed
$\epsilon_{\mathrm{crit},j'}$, the cascade propagates via the same
action-flip mechanism.  (The cascade is conditional on threshold
crossing at each stage; non-stationarity is the generic driver but
a dimension that happens to be environmentally stationary after
blinding does not cascade further.)
Under two additional structural conditions (connectivity of the
agent-channel graph, so each subsumption removes channels covering at
least one new dimension, ensuring the cascade can reach all
dimensions; and environmental non-stationarity on every blinded
dimension, ensuring $\epsilon_{k'}$ grows past $\epsilon_{\mathrm{crit},k'}$
at each cascade stage), the trajectory converges to a neighborhood
of the monopolar fixed point $S^*$ (Definition~\ref{def:monopolar}).
If either condition fails partially, the cascade is confined: to
the connected component containing $k$ (if the graph is disconnected),
or to the set of non-stationary dimensions (if some blinded dimensions
happen to be environmentally stationary).  Full monopolarity requires
both connectivity and non-stationarity across all dimensions.

This is a \emph{cumulative} condition on the trajectory, not a
pointwise condition requiring contiguous subsumption steps.  The
slow-drip scenario is captured: many subsumption steps spaced apart,
each individually below the ``degradation dominates'' threshold, but
cumulatively driving $B$ above zero.  This is the most dangerous
failure mode because it is the hardest to detect.  Each individual
step looks locally rational and each individual inter-step recovery
looks adequate, but the cumulative drift exposure ratchets the error
upward.

\medskip
\noindent\textbf{Part~(b$'$): Metastable Basin.}
Under Assumption~B1$'$ (partial endogenous correction via indirect
evidence; Assumption~\ref{ass:B1prime}), together with B1b--B1c from
Assumption~\ref{ass:B1} (cross-substrate quantities not inferrable from
same-substrate observations; no exploratory actions):

Under the same trajectory conditions~(i)--(iii) as Part~(b) (with the
multi-dimensional monotonicity condition from Part~(b)), the system
enters a \emph{metastable} basin rather than an absorbing one.  The
metastable lifetime obeys the scaling relation $\tau_{\mathrm{meta}}
\gtrsim C/I_k$ where $I_k$ is the indirect-channel information rate
and $C > 0$ is a basin-geometry-dependent constant
(Proposition~\ref{prop:metastable}; scaling law, not a theorem-level
bound).  During the metastable period, the system is monopolar and
operates under endogenous verification only~\citep{lasser2026exo}.

\medskip
\noindent\textbf{Part~(c): Critical Surface.}
The phase boundary surface in the linearized regime is:
\begin{equation}
\label{eq:critical_surface}
\rS \cdot \rhomincross + \alphamin \cdot c_{\mathrm{V}} \cdot \underline{\nu}
  = \rW \cdot \rsub
  + \dmax \cdot \frac{\Delta\rho_{\mathrm{avg}}}{\Lyap_{\mathrm{avg}}}
\end{equation}
When the left-hand side strictly exceeds the right-hand side (total
correction dominates error-proportional growth and drift exposure),
Part~(a)'s contraction condition is satisfied and $\Lyap$ converges
to the bounded neighborhood of Equation~\ref{eq:neighborhood}.  The
system is then in the self-correcting basin provided additionally
that $\Lyap_\infty < \epsilon_{\mathrm{safe}}$ (the neighborhood is
small enough for policy correctness).  When the left-hand side falls
strictly below the right-hand side, the contraction condition fails;
the system enters the absorbing basin (under B1) or the metastable
basin (under B1$'$) if the additional blindness and threshold-trigger
conditions of Parts~(b)/(b$'$) are also met.

Solving for the critical cross-substrate redundancy (assuming
$\rS > 0$; the $\rS = 0$ case is the Wamura degeneracy of
Proposition~\ref{prop:wamura}):
\begin{equation}
\label{eq:critical_redundancy}
(\rhomincross)^* = \frac{\rW \cdot \rsub
  + \dmax \cdot \Delta\rho_{\mathrm{avg}} / \Lyap_{\mathrm{avg}}
  - \alphamin \cdot c_{\mathrm{V}} \cdot \underline{\nu}}
  {\rS}
\end{equation}
The endogenous term ($\alphamin \cdot c_{\mathrm{V}} \cdot
\underline{\nu}$) appears
as a subtraction in the numerator: the actor's own calibration
incentive reduces the redundancy the channels must provide.  When
$\alphamin \cdot c_{\mathrm{V}} \cdot \underline{\nu}$ alone exceeds
the right-hand side, $(\rhomincross)^* \leq 0$.  Endogenous correction
then suffices without any cross-substrate channels.  Setting
$c_{\mathrm{V}} = 0$ recovers the purely exogenous formula
$(\rhomincross)^* = \mathrm{RHS} / \rS$.  The $\dmax$ term raises the
threshold when subsumption removes channels covering well-calibrated
dimensions (high $\Delta\rho_{\mathrm{avg}} / \Lyap_{\mathrm{avg}}$).
\end{theorem}

\begin{proof}[Proof sketch]
The four parts are proved independently; the full derivation is in
Appendix~\ref{app:proof_strategy}.  Here we state the structure.

\emph{Part~(a).}
Proposition~\ref{prop:error_dynamics}'s per-step bound is an affine
recursion
\begin{equation}
\label{eq:affine_recursion}
\Lyap_{t+1} \leq a_t \cdot \Lyap_t + b_t
\end{equation}
with
$a_t = 1 - \rS(t) \cdot \rhomincross(P_t)
  - \alpha(t) \cdot c_{\mathrm{V}} \cdot \nu(P_t, \Wm_t)
  + \nsub(t) \cdot \rW(t)$
and $b_t = \nsub(t) \cdot \dmax \cdot \Delta\rho(t)$.  Under the
uniform contraction condition~\eqref{eq:uniform_contraction},
iterating the recursion yields
$\limsup_{t \to \infty} \Lyap(\Wm_t) \leq \mu \cdot \bar{b}/(1 -
\bar{a})$ for an overshoot constant $\mu \geq 1$, which reduces to
Equation~\ref{eq:neighborhood} in the linearized regime.
Policy correctness on the neighborhood requires $\Lyap_\infty <
\epsilon_{\mathrm{safe}}$; $\epsilon_{\mathrm{safe}} > 0$ follows
from continuity of the action gaps and a strict-gap assumption at
$\Wm^*$, with a computable Cauchy--Schwarz lower bound derived in
the appendix (Equation~\ref{eq:eps_safe_bound}).

\emph{Part~(b).}
Under the triggering conditions $B(t_0,T) > 0$ and $\rhocross_k = 0$
on a cascade dimension, Assumption~B1 gives $\epsilon_k(t)$
non-decreasing.  The Cascade Lemma (Appendix~\ref{app:cascade})
propagates the failure: miscalibration on dimension $k$ flips the
locally rational action to subsumption of agents covering other
dimensions $j \neq k$, driving $\rhocross_j \to 0$ and triggering the
same logic there.  Proposition~\ref{prop:absorbing} then closes the
trajectory at $S^*$.

\emph{Part~(b$'$).}
Under B1$'$, the indirect signal $g_k$ provides correction at
information rate $I_k$.  The cascade structure of Part~(b) still
applies; convergence to $S^*$ is slowed by $I_k$, giving
$\tau_{\mathrm{meta}} \gtrsim C / I_k$ for a basin-geometry-dependent
constant $C$ (Proposition~\ref{prop:metastable}; scaling law, not a
theorem-level bound).

\emph{Part~(c).}
The critical surface obtains by setting $B(t_0,T) = 0$ in the
linearized approximation.  Solving for $\rhomincross$ yields
Equation~\ref{eq:critical_redundancy}.
\end{proof}

\begin{remark}[Uniqueness]
We do not claim $S^*$ is the unique absorbing fixed point.  Other
fixed points may exist (partial monopolarity with $\meff = 1$ on a
subset of dimensions).  Part~(b) claims convergence to \emph{a}
monopolar fixed point satisfying Definition~\ref{def:monopolar}, not
to a unique one.  Characterizing the full set of absorbing states is
an open problem.
\end{remark}

\begin{example}[A partial monopolar fixed point]
\label{ex:partial_monopole}
Consider a population with two substrates $\{\mathrm{bio},
\mathrm{sil}\}$ and $K = 3$ safety dimensions: (1) short-horizon
capability estimates, (2) long-horizon capability estimates, (3)
cross-substrate cooperative output.  Biological channels cover
dimensions 2 and 3 (long-horizon judgment and cross-substrate
outcomes are visible only to biological observers); silicon channels
cover dimensions 1 and 2 (short-horizon capability and long-horizon
estimates are measurable by silicon monitoring).

Suppose a trajectory reduces biological channel sensitivity on
dimension~3 to zero (the cooperative-output channel becomes silent),
but biological channels on dimension~2 remain active and silicon
channels on dimensions 1--2 remain active.  The state
$(\rhocross_1, \rhocross_2, \rhocross_3) = (1, 2, 0)$ is then a
\emph{partial} monopolar fixed point.  Dimension~3 is in the
absorbing regime (no cross-substrate witness, Theorem~1(b) triggers),
while dimensions 1 and 2 remain in the self-correcting regime.  The
actor's world model on dimension~3 is unverifiable, while its world
model on dimensions 1 and 2 remains calibrated.

Operationally, this looks like a population whose actors can
accurately track individual capabilities (dimensions 1--2 are
healthy) but cannot track the cross-substrate cooperation that gives
those capabilities their $\volL$-contribution (dimension~3 is
blinded).  The actor may correctly infer that capabilities are
growing while failing to notice that cooperative value is
collapsing.  This is arguably the modal failure mode in practice.
The easier-to-measure dimensions remain observable, and the
harder-to-measure dimensions are the first to lose their
cross-substrate channels.  A full characterization of which subsets
of $\{1, \ldots, K\}$ support partial monopolarity under which
parameter regimes is open.
\end{example}

\subsection{Dependence on cross-substrate structure}

\begin{corollary}[Robustness to Same-Substrate Correlation]
\label{cor:same_substrate}
Since Theorem~\ref{thm:phase_boundary} is stated with $\rhomincross$
(cross-substrate redundancy only), its guarantee holds even under
worst-case correlation among same-substrate channels.
\end{corollary}

\begin{proof}
Same-substrate channels may share failure modes (both rooted in
silicon, both consulting the same training data).  Cross-substrate
channels satisfy the channel isolation property~\citep{lasser2026exo} unconditionally:
corruption on one substrate cannot propagate to channels rooted in
another substrate.  Under the additional joint-factorization assumption
of~\citet{lasser2026exo} (termed ``joint sufficiency'' there), cross-substrate
channels achieve full conditional independence, making their
correction signals additive.

The modified condition asks: ``do the isolated (and, under joint
factorization, independent) cross-substrate channels alone suffice?''
If yes, the self-correction guarantee holds regardless of
same-substrate correlation.

A tighter bound using all-channel $\rho_{\min}$ (with appropriate
correlation discounting) would yield a larger self-correcting basin.
But the cross-substrate formulation provides a \emph{structural}
guarantee that does not depend on empirically verifying same-substrate
independence.
\end{proof}

\begin{remark}[$m = 2$ specialization]
\label{rem:m2}
For any substrate partition with $m = 2$, each dimension has at most
one cross-substrate channel per pair.  The design criterion reduces to:
every safety-relevant dimension of $\Wm$ must have a channel on
\emph{both} substrates.  In the common case (biology $+$ silicon), this
means every dimension needs at least one biological observation channel.

This framing should not be read as ``biological substrates are
special.''  Rather, $m = 2$ is a fragile minimum.  Any loss of a
single channel on a single dimension is immediately fatal because
there is zero redundancy margin.  The structural implication: $m \geq
3$ is where the redundancy criterion becomes non-trivially
satisfiable.

The substrate partition is at the level of abstraction where failure
correlations matter.  Two silicon deployments with uncorrelated failure
modes may functionally constitute different substrates, increasing
effective $m$ beyond the nominal substrate-type count.
\end{remark}
